\newcommand{\CaseID}{caldor\_2021}
\newcommand{\CaseTitle}{2021 Caldor Fire Reconstruction against GOFER Hourly Perimeters}
\newcommand{\FireName}{Caldor Fire}
\newcommand{\FireYear}{2021}
\newcommand{\EventStart}{15 August 2021 at 01:00~UTC}
\newcommand{\EventEnd}{10 September 2021 at 03:00~UTC}
\newcommand{\DurationHours}{626}
\newcommand{\CRSName}{UTM zone 10N (EPSG:32610)}
\newcommand{\FuelVersionText}{LANDFIRE release 2.2.0}
\newcommand{\FireSummary}{GOFER records the Caldor Fire from 15 August 2021 at 01:00~UTC to 10 September 2021 at 03:00~UTC with 620 hourly perimeters, a final perimeter area of 994~km$^2$ and an official size of 221,835~ac, and 1 ignition point(s).}
\newcommand{\IgnitionMoveNote}{no point required moving}
\newcommand{\OffshoreNote}{}
\newcommand{\CalibrationNote}{No model parameter was tuned against the GOFER observation of this fire; the spotting configuration was revised after comparisons for the Kincade and CZU Lightning Complex fires.}
\newcommand{\ConfigurationText}{The domain is the 30-m landscape grid in UTM zone 10N (EPSG:32610).  The simulation starts at the first GOFER ignition point and runs 620~h, to the last GOFER perimeter, with a single point ignition.  Weather and fuel moisture are read hourly and interpolated in time.  Surface spread follows the Rothermel model with the Scott and Burgan fuel models.  Crown fire uses the ELMFIRE default treatment: initiation when surface fireline intensity exceeds the critical intensity of Van Wagner \cite{vanwagner1977}, and active-crown spread rate from Cruz et al. \cite{cruz2005}, capped at 250~ft~min$^{-1}$, with a constant foliar moisture of 90\%.  Firebrand spotting is enabled with every other spotting setting at its ELMFIRE default, so firebrands are generated only by crown-fire cells.  The level-set equation is advanced with an adaptive timestep targeting a Courant--Friedrichs--Lewy (CFL) number of 0.2, starting at 1~s and limited to 600~s.  The run is a single deterministic member with a fixed random seed; arrival time, spread rate, fireline intensity and fire-size statistics are recorded at the stop time.}
\newcommand{\WeatherOffsetText}{6~h after the GOES ignition hour; all results are reported in hours since the GOES ignition hour.}
\newcommand{\RevisionText}{The simulation results reported here were produced with ELMFIRE~1.1 at source revision \texttt{5ea7460d514b}, as recorded for the CZU Lightning Complex runs and assumed for the other results produced on the same system.}
\newcommand{\CaseInterpretation}{}
\newcommand{\ElmfireRevision}{5ea7460d514b}
